%% Copernicus Publications Manuscript Preparation Template for LaTeX Submissions
%% ---------------------------------
%% This template should be used for copernicus.cls
%% The class file and some style files are bundled in the Copernicus Latex Package, which can be downloaded from the different journal webpages.
%% For further assistance please contact Copernicus Publications at: production@copernicus.org
%% https://publications.copernicus.org/for_authors/manuscript_preparation.html


%% Please use the following documentclass and journal abbreviations for discussion papers and final revised papers.

%% 2-column papers and discussion papers
\documentclass[journal abbreviation, manuscript]{copernicus}



%% Journal abbreviations (please use the same for discussion papers and final revised papers)


% Advances in Geosciences (adgeo)
% Advances in Radio Science (ars)
% Advances in Science and Research (asr)
% Advances in Statistical Climatology, Meteorology and Oceanography (ascmo)
% Annales Geophysicae (angeo)
% Archives Animal Breeding (aab)
% ASTRA Proceedings (ap)
% Atmospheric Chemistry and Physics (acp)
% Atmospheric Measurement Techniques (amt)
% Biogeosciences (bg)
% Climate of the Past (cp)
% DEUQUA Special Publications (deuquasp)
% Drinking Water Engineering and Science (dwes)
% Earth Surface Dynamics (esurf)
% Earth System Dynamics (esd)
% Earth System Science Data (essd)
% E&G Quaternary Science Journal (egqsj)
% Fossil Record (fr)
% Geochronology (gchron)
% Geographica Helvetica (gh)
% Geoscience Communication (gc)
% Geoscientific Instrumentation, Methods and Data Systems (gi)
% Geoscientific Model Development (gmd)
% History of Geo- and Space Sciences (hgss)
% Hydrology and Earth System Sciences (hess)
% Journal of Micropalaeontology (jm)
% Journal of Sensors and Sensor Systems (jsss)
% Mechanical Sciences (ms)
% Natural Hazards and Earth System Sciences (nhess)
% Nonlinear Processes in Geophysics (npg)
% Ocean Science (os)
% Primate Biology (pb)
% Proceedings of the International Association of Hydrological Sciences (piahs)
% Scientific Drilling (sd)
% SOIL (soil)
% Solid Earth (se)
% The Cryosphere (tc)
% Web Ecology (we)
% Wind Energy Science (wes)


%% \usepackage commands included in the copernicus.cls:
%\usepackage[german, english]{babel}
%\usepackage{tabularx}
%\usepackage{cancel}
%\usepackage{multirow}
%\usepackage{supertabular}
%\usepackage{algorithmic}
%\usepackage{algorithm}
%\usepackage{amsthm}
%\usepackage{float}
%\usepackage{subfig}
%\usepackage{rotating}


\begin{document}

\title{Synergistic radar and sub-millimeter radiometer retrievals of ice clouds}

% \Author[affil]{given_name}{surname}

\Author[1]{Simon}{Pfreundschuh}
\Author[1]{Patrick}{Eriksson}
\Author[2]{Stefan A.}{Buehler}
\Author[2]{Manfred}{Brath}
\Author[4]{David}{Duncan}
\Author[3]{Richard}{Larsson}
\Author[2]{Robin}{Eklund}

\affil[1]{Department of Space, Earth and Environment, Chalmers University of Technology, 41296 Gothenburg, Sweden}
\affil[2]{Meteorologisches Institut, Fachbereich Geowissenschaften, Centrum für Erdsystem und Nachhaltigkeitsforschung (CEN), Universität Hamburg, Bundesstraße 55, 20146 Hamburg, Germany}
\affil[3]{Max Planck Institute for Solar System Research, Justus-von-Liebig-Weg 3, 37077 Göttingen, Germany}
\affil[4]{ECMWF, Shinfield Park, Reading RG2 9AX, United Kingdom}
%% The [] brackets identify the author with the corresponding affiliation. 1, 2, 3, etc. should be inserted.



\runningtitle{Retrieving frozen hydrometeors from combined radar and sub-millimter observations}
\runningauthor{Simon Pfreundschuh}
\correspondence{Simon Pfreundschuh (simon.pfreundschuh@chalmers.se)}



\received{}
\pubdiscuss{} %% only important for two-stage journals
\revised{}
\accepted{}
\published{}

%% These dates will be inserted by Copernicus Publications during the typesetting process.


\firstpage{1}

\maketitle



\begin{abstract}
  The upcoming Ice Cloud Imager (ICI) instrument, to be launched on board the
  second generation of European operational meteorological satellites (MetOpSG),
  will be the first microwave imager to provide sub-millimeter observations of
  clouds. To investigate what can be learned from these newly available
  frequencies, this study explores the potential of combining observations from
  a sub-millimeter wave radiometer with those from a $94 \unit{GHz}$ cloud-radar
  for the retrieval of frozen hydrometeors. Starting from a simplified numerical
  experiment, it is shown that the complementary information content in the
  radar and radiometer observations helps to better constrain the particle size
  distribution of ice particles in the atmosphere. The feasibility of the
  combined retrieval is demonstrated by applying a one-dimensional, variational
  cloud-retrieval algorithm to simulated observations from a high-resolution
  atmospheric model. Comparison of the retrieval results with passive- and
  radar-only versions of the retrieval algorithm confirm that the synergistic
  information content from the active and passive observations helps to better
  constrain the microphysical properties of frozen hydrometeors. In addition to
  that, the synergistic retrieval shows potential to detect and retrieve liquid
  clouds. The results of this study clearly demonstrate the potential of
  combining radar and sub-millimeter radiometer observations for the study of
  ice clouds. As observational concept, the combined retrievals point out a
  potential way forward towards reducing uncertainties in the observational
  record of ice in the atmosphere.
\end{abstract}


\introduction  %% \introduction[modified heading if necessary]

Ice clouds play an important role in many weather- and climate-related
processes in the atmosphere. They interact with incoming and outgoing
radiation and thus influence the Earth's energy budget. Furthermore, as part
of the global hydrological cycle and due to their relation to the dynamics of
the atmosphere, observations of ice clouds provide important information to
constrain the state of the atmosphere in forecasting models \citep{geer} as well as
to validate predictions from climate models \citep{waliser09}.

Despite the importance of observations of ice clouds for climate and weather
prediction, todays global earth observation system cannot provide
accurate information on the distribution of ice in the atmosphere
\citep{eliasson11,duncan18a}. What makes the remote sensing of ice from space
so difficult is the large variability in sizes and concentrations in which
ice particles occur in the atmosphere. The spectrum of ice crystal sizes,
that ranges from micro- to millimeter scales, can only be partially resolved
by currently available space-borne sensors.

The sensitivity of a remote sensing system to ice particles of a given size is
determined mainly by its observing frequencies. The interaction of radiation
with ice particles is strongest for sizes roughly equal to the wavelength,
$\lambda$, of the radiation. For particles with sizes smaller than $\lambda$,
the sensitivity decreases so rapidly that they are practically invisible to the
sensor. Particles larger than the wavelength $\lambda$ interact sufficiently
strongly with the radiation, that their presence quickly leads to saturation of
the signal and thus loss of sensitivity further down the line of sight.

The observing frequencies that are currently available for measuring ice from
space are limited to the microwave, infrared and optical domain. Infrared and
optical sensors, on the one hand, provide sensitivity to small ice particles but
cannot sense significant parts of the ice mass of thicker clouds due to
saturation of the signal. Microwave observations, on the other hand, provide
sensitivity throughout the whole atmospheric column but are insensitive to small
ice particles. Although radars and lidars generally provide higher sensitivity
than their passive counterparts, ultimately their sensitivity is limited by the
same principles.

To narrow this sensitivity gap, the upcoming Ice Cloud Imager (ICI) will extend
the microwave frequencies available for studying clouds with channels at $243,
325, 448$ and $664\unit{GHz}$. This extension of the smallest currently
available microwave wavelength from $1.6\unit{mm}$ at $183\unit{GHz}$ down to
the sub-millimeter domain ($0.45\unit{mm}$ at $664\unit{GHz}$) will
significantly improve the size-sensitivity of space-borne microwave observations
of clouds.

With the advent of space-borne sub-millimeter radiometry of clouds on the
horizon, a natural question to ask is what can be learned from this new
frequency domain about  ice in the atmosphere. For passive
radiometer observations alone, this has been studied in detail for
column-integrated ice mass \citep{jimenez07, wang16, brath18a} as well as for
the vertical distribution of ice in the atmosphere \citep{wang16, grutzun18}.
However, the concept of combining sub-millimeter observations with active
observations from a cloud radar has remained mostly unexplored since it was
first proposed by \cite{evans05}.

This study revisits this concept and explores its potential for retrieving
vertically-resolved distributions of frozen hydrometeors in the atmospheric
column. The work presented here was performed as part of a larger study funded
by the European Space Agency, which aimed to evaluate the concept of a future
radar mission to fly in constellation with ICI on board the satellites of the
MetOp-Second Generation Program. Inspired by the success of the Global
Precipitation Measurement (GPM, \cite{hou14}), the approach of the mission is to
perform vertically-resolved, high-accuracy retrievals of hydrometeors from the
co-located active and passive observations at the swath center of the passive
imager. The combined retrieval results are then used to constrain passive-only
profile retrievals with the aim of extending the profiling capabilities of the
radar to the wide swath of the passive imager.

The aim of this study is to explore the fundamental synergies between radar and
passive sub-millimeter observations. In particular, it aims to address the
question to what extent the combined active and passive obervations can
constrain the microphysics of ice particles in the atmosphere. Starting from
a simplified numerical experiment, the complementarity of the information content
of the active and passive observations is demonstrated. Following this, a
synergistic, variational cloud-retrieval algorithm is presented. The algorithm
is applied to simulated observations of cloud scenes from a high-resolution
atmospheric model and used to further explore the synergies between the active
and passive observations.

Following this introduction, Section \ref{sec:methods_and_data} introduces
the test data, sensor configuration and the developed retrieval algorithm
on which the study is based. This is followed by the presentation of the
experimental results on the information content of the combined observations
and the retrieval results of the joint retrieval on selected test scenes
in Section \ref{sec:results}. The article closes with a discussion of
the results in Section \ref{sec:discussion} and conclusions in section
\ref{sec:conclusions}.


\section{Methods and data}
\label{sec:methods_and_data}

For the performance assessment of the synergistic cloud retrieval, radiative
transfer simulations of cloud scenes from a high-resolution climate model are
used. This section introduces the selected reference cloud scenes, the sensor
configuration and the modeling assumptions for the simulation of the synthetic
observations. Furthermore, the theoretical formulation of the combined
cloud-retrieval algorithm is presented.

\subsection{Reference cloud scenes}

The cloud scenes that are used to generate the synthetic cloud observations are
taken from the Global Environmental Multiscale Model (GEM, \cite{cote98}). For
this study, we restrict ourselves to two designated, two-dimensional test
scenes, that are displayed in Fig.~\ref{fig:overview}. The test scenes have a
horizontal resolution of $1\ \unit{km}$ and both extend over $800\ \unit{km}$.
The scenes where chosen with the aim of covering a large range of cloud
structures and compositions so as to ensure a realistic assessment of the
retrieval. The first test scene, shown in panel (a), is located in the tropical
Pacific and contains a convective storm system in the right half of the scene
and its anvil that extends into the left half of the scene. The second scene,
shown in panel (b), is located in the North Atlantic and contains an  ice
cloud in the first quarter and a low-level, mixed-phase cloud in the remainder
of the scene.

\begin{figure}[h!]
\centering
\includegraphics[width = 0.8\textwidth]{../plots/scene_overview.png}
\caption{The distribution of total hydrometeor mass content in the two
cloud scenes used to test the retrieval. Colored lines show the
 $m = 10^{-5} \unit{kg/m^3}$ contour for different
 hydrometeor species.}
\label{fig:overview}
\end{figure}


The GEM model, from which the scenes are taken, uses six types of hydrometeors
to represent clouds and precipitation \citep{milbrandtyau05}: Two classes of
liquid hydrometeors (rain and liquid cloud) and four of frozen hydrometeors
(cloud ice, snow, hail and graupel). The particle size distribution (PSD) of
each hydrometeor type is parametrized by its particle number concentration and
mass density. The full particle size distribution can be prognosed from the two
moments using a species-dependent parametrization and mass-size relationship.
The parameters of the mass-size relationship are given in
Tab.~\ref{tab:species_parameters}.


\begin{table}
  \centering
  \begin{tabular}{l|c|c|c|c}
    Hydrometeor species & Particle shape & ID & $\alpha$ & $\beta$ \\
    \hline
    Cloud ice    & GemCloudIce  & 31 & 440   & 3 \\
    Snow         & GemSnow      & 32 & 52.4  & 3 \\
    Graupel      & GemGraupel   & 33 & 209.4 & 3 \\
    Hail         & GemHail      & 34 & 471.2 & 3 \\
    Rain         & LiquidSphere & 25 & 523.6 & 3 \\
    Liquid cloud & LiquidSphere & 25 & 471.2 & 3 \\
  \end{tabular}
  \caption{Particle-model names, IDs and parameters $\alpha, \beta$ of the
    mass-size relationships $m = \alpha D_\text{max}^\beta$ used for simulating
    observations from the GEM model scenes. The ID column contains the particle
    shape identifier of the particle model in the ARTS scattering database
    \citep{eriksson18}.}
  \label{tab:species_parameters}
\end{table}

Examples of full particle size distributions of frozen hydrometeors are
displayed in Fig.\ref{fig:gem_psds}. The four panels display the prognosed
particle size distributions for the four frozen hydrometeor types together with
renderings of the particle shapes used in the forward simulations. As these
plots show, the assumed particle size distributions across different ice species
vary mostly in their horizontal and vertical scaling, whereas the function shape
shows less variability. Already here an important characteristic of the model
can be seen, that will help to better understand the retrieval results: Cloud
ice in the model is characterized by high particle number densities and small
particle sizes, while snow exhibits lower number concentrations and larger
particles.


\begin{figure}[h!]
\centering
\includegraphics[width = \textwidth]{../plots/gem_hydrometeors.png}
\caption{Realizations of particle size distributions from the cloud scenes used in
this study. Shown are the PSDs corresponding to 100 randomly chosen grid points at
which the mass concentration was higher than $10^{-5}$. Line color encodes the
corresponding mass density.}
\label{fig:gem_psds}
\end{figure}


\subsection{Simulated cloud observations}

For the generation of the synthetic cloud observations that are used to test the
retrieval, the full GEM microphysics scheme is used. This means that the forward
simulations use the same six hydrometeor classes as the GEM model. In addition to
that, temperature, pressure, altitude and humidity of the simulated atmosphere
are taken from the model scenes. Profiles of oxygen and nitrogen are taken from
the Fascod climatology for tropical climates.

\subsubsection{Sensor configuration}
\label{sec:sensors}

The assumed sensor configuration for the combined retrieval consists of the
MWI and ICI radiometers and a W-band cloud radar similar to CloudSat CPR
\citep{stephens02,tanelli08}. The choice of a CloudSat CPR-type radar is based
mainly on the success and longevity of the CloudSat mission, that clearly
demonstrates the maturity and robustness of the technology.

A number of simplifications are applied for the generation of the synthetic
cloud observations: Firstly, the beams of all three sensors are modeled as
perfectly coincident pencil beams. Secondly, for each vertical profile from the
model scenes a synthetic observation is generated by simulating a
one-dimensional, plane-parallel atmosphere, the properties of which are taken
from the corresponding model profile. It follows from these modeling decisions
that the atmosphere is assumed to be homogeneous across the beams of the active
and passive sensors and that they all sense the same atmospheric volume. For
space-borne observations, this would certainly not be the case and will incur a
forward modeling error that is not accounted for in this study. The focus of
this study are the fundamental synergies between the active and passive
observations, which is why this error is neglected here. It will be considered
in a follow-up study that evaluates the synergistic retrievals in the context of
a potential future space mission.

Observations from the ICI radiometer are simulated by performing a single,
non-polarized radiative transfer simulation located at the centers of the side
bands of each channel and averaging the resulting brightness temperatures. For
channels with multiple polarizations, only a single simulation is performed.
To compensate for this, the noise of the corresponding channel is reduced by a
factor of $\sqrt{2}$. The simulated ICI channels and assumed noise levels are
presented in  of Table \ref{tab:channels}. The off-nadir viewing angle of ICI
is assumed to be $48\ \unit{^\circ}$.

Observations from the MWI radiometer are simulated in a similar manner as for
ICI. However, from MWI only channels with frequencies larger than or equal to
$89 \unit{GHz}$ are used. The reason for this is that the footprints of the
channels with frequencies lower than $89\ \unit{GHz}$ will have full-width at
half maximum of $50 \unit{km}$ compared to only $10 \unit{km}$ for the
high-frequency channels. It is therefore assumed that, due to the very little
expected overlap of the footprints of the low-frequency channels with a radar,
they would not be beneficial for a synergistic retrieval and are therefore
disregarded here. The included MWI channels are listed in Table
\ref{tab:channels}.

\begin{table}[hbpt]
\caption{Channels of the MWI and ICI radiometers used in the retrieval.}
\label{tab:channels}
    \begin{tabular}{c|r|r}
    \multicolumn{3}{c}{MWI}\\
    Channel & Freq. [GHz] & Noise [K]\\
    \hline
    MWI-8  & $89$              & $1.1$ \\
    MWI-9  & $118.75 \pm 3.2$  & $1.3$ \\
    MWI-10 & $\pm 2.1$         & $1.3$ \\
    MWI-11 & $\pm 1.4$         & $1.3$ \\
    MWI-12 & $\pm 1.2$         & $1.3$ \\
    MWI-13 & $165.5 \pm 0.75$  & $1.3$ \\
    MWI-14 & $183.31 \pm 7.0$  & $1.2$ \\
    MWI-15 & $ \pm 6.1$        & $1.2$ \\
    MWI-16 & $ \pm 4.9$        & $1.2$ \\
    MWI-17 & $ \pm 3.4$        & $1.2$ \\
    MWI-18 & $ \pm 2.0$        & $1.3$ \\
    \end{tabular}%
    \hspace{1cm}%
    \begin{tabular}{c|r|r}
    \multicolumn{3}{c}{ICI}\\
    Channel & Freq. [GHz] & Noise [K] \\
    \hline
    ICI-1  & $183.31 \pm 7.0$ & $0.8$\\
    ICI-2  & $       \pm 3.4$ & $0.8$\\
    ICI-3  & $       \pm 2.0$ & $0.8$\\
    ICI-4  & $243    \pm 2.5$ & $\frac{1}{\sqrt{2}} \cdot 0.7$\\
    ICI-5  & $325.15 \pm 9.5$ & $1.2$\\
    ICI-6  & $       \pm 3.5$ & $1.3$\\
    ICI-7  & $       \pm 1.5$ & $1.5$\\
    ICI-8  & $448    \pm 7.2$ & $1.4$\\
    ICI-9  & $       \pm 3.0$ & $1.6$\\
    ICI-10 & $       \pm 1.4$ & $2.0$\\
    ICI-11 & $664    \pm 4.2$ & $\frac{1}{\sqrt{2}} \cdot 1.6$\\
    \end{tabular}
\end{table}

The frequency of the the cloud radar is chosen to be $94\unit{GHz}$ similar to
CloudSat CPR. The vertical resolution of the radar observations is assumed to be
$500\unit{m}$ ranging from $0.5$ to $20\unit{km}$ in altitude. The minimum
sensitivity is set to be $-30\unit{dBZ}$ and the noise at each range gate is
modeled to be independent with standard deviation $0.5\unit{dBZ}$. As mentioned above,
the same incidence angle as for the passive radiometers is assumed also for the
radar. In reality, this could be achieved by remapping the radar observations to
the lines of sights of the passive beams.

\subsubsection{Radiative transfer simulations}
\label{sec:orge741b86}

All simulations presented in this study were performed using Version 2.3.1245 of
the Atmospheric Radiative Transfer Simulator (ARTS, \cite{arts18}). Radar
reflectivities are computed using the ARTS' built-in single-scattering radar
solver. For the simulation of passive radiances, a hybrid solver is used that
combines the DISORT scattering solver with ARTS built-in line-by-line solver.
The hybrid solver has been added to ARTS specifically for this study and
provides approximate Jacobians, which are required to perform variational
retrievals of hydrometeors. All simulations are performed assuming an
ocean surface with emissivities calculated using the Tool to Estimate
Sea‐Surface Emissivity from Microwaves to sub‐Millimeter waves (TESSEM,
\cite{prigent16}). Polarization is neglected in all simulations performed in
this study. Particle scattering data are taken from the ARTS scattering data base
\citep{eriksson18}.

\subsection{Retrieval algorithm}
\label{sec:orgb528563}

A one-dimensional, variational cloud retrieval algorithm is proposed to
retrieve distributions of frozen hydrometeors from the combined active and
passive observations. The algorithm uses the optimal estimation formalism
developed by \cite{rodgers00}. The retrieved state $\mathbf{x} \in
   \mathbb{R}^n$ is determined by fitting a forward model $\mathbf{F} : \mathbb{R}^n
   \rightarrow \mathbb{R}^m$ to a set of observations $\mathbf{y} \in
   \mathbb{R}^m$. The best fit is determined by minimizing a cost function of
the form

\begin{align}
l(\mathbf{x}, \mathbf{y}) \propto
 \left(\mathbf{F}(\mathbf{x}) - \mathbf{y} \right )^T
  \mathbf{S}_e^{-1} 
  \left ( \mathbf{F}(\mathbf{x}) - \mathbf{y} \right)
+ \left ( \mathbf{x} - \mathbf{x}_a \right )^T
 \mathbf{S}^{-1}_a 
 \left ( \mathbf{x} - \mathbf{x}_a \right ).
\end{align}

The cost function $l(\mathbf{x}, \mathbf{y})$ corresponds to the negative
log-likelihood of the a posteriori distribution of the state $\mathbf{x}$
under the assumptions of Gaussian a priori distribution with mean
$\mathbf{x}_a$ and covariance matrix $\mathbf{S}_a$ as well as zero-mean Gaussian
measurement error with covariance
matrix $\mathbf{S}_e$.

\subsubsection{Measurement space}
\label{sec:orge7dc286}

The input for the retrieval algorithm is the combined observation vector
$\mathbf{y}$, which consists of the concatenated single-instrument observation
vectors. Measurement errors are assumed to be independently Gaussian-distributed
with standard deviations according to the noise characteristics given in Section
\ref{sec:sensors}. 

\subsubsection{State space}
\label{sec:orgae5ee4e}

The proposed retrieval algorithm retrieves distributions of one frozen
and one liquid hydrometeor species in the atmospheric column together with
profiles of atmospheric humidity and liquid-cloud mass density.

The distributions of hydrometeors in the atmospheric column are represented
using the normalized particle size distribution formalism proposed by
\cite{delanoe05}. The PSD of a hydrometeor species at a given grid point is
represented by a vertical and a horizontal scaling parameter, the normalized
number density $N_0^*$ and the mass-weighted mean diamter $D_m$. The shape of
the PSD, described by the parameters $\alpha$ and $\beta$, is assumed fixed for
each hydrometeor species. For frozen hydrometeors, $\alpha$ and $\beta$ are
chosen identical to version 3 of the DARDAR-CLOUD product \citep{cazenave18}.
For liquid hydrometeors, the shape parameters are chosen so that they are
equivalent to the shape used by the GEM model for rain drops. Figure
\ref{fig:psds_retrieval} displays the resulting PSD shapes and corresponding
parameter values.

\begin{figure}
\centering
\includegraphics[width = 0.5\textwidth]{../plots/psds_retrieval}
\caption{PSD parametrizations for frozen and liquid hydrometeors
 used in the cloud retrieval.}
\label{fig:psds_retrieval}
\end{figure}

The temperature-dependent a priori profile for $N_0^*$ of frozen
hydrometeors is determined using the relation from \cite{delanoe14}

\begin{align}
N_0^* &= \exp (-0.076586 \cdot t + 17.948),
\end{align}

where $t$ is in $\unit{^\circ C}$. The a priori profile for $D_m$ for frozen
hydrometeors is chosen so that the a priori mass density is equal to
$10^{-6}\ \unit{kg/m^3}$. For liquid hydrometeors, a
fixed value for $N_0^*$ of $10^6\ \unit{m^4}$ is assumed and the
a priori profile for $D_m$ is determined similarly as for frozen
hydrometeors. Values of the mass-weighted mean diameter $D_m$ for both
hydrometeor species are retrieved in linear space, whereas the normalized
number concentration parameter $N_0^*$ is retrieved in $\text{log}_{10}$
space. As additional constraints, the retrieval of frozen hydrometeors is
restricted to the region between the $0\unit{^\circ C}$-isotherm and
the tropopause, whereas the retrieval of liquid hydrometeors is restricted 
to below the $0\unit{^\circ C}$-isotherm.

Humidity in the atmospheric column is retrieved in units of relative humidity.
However, instead of retrieving relative humidity directly, an inverse hyperbolic
tangens transformation is applied to the relative humidity profile $\mathbf{\phi}$:

\begin{align}
x = \text{arctanh}(\frac{2 \mathbf{\phi}}{1.2} - 1.0)
\end{align}

The transformation restricts the retrieved relative humidity values to
the range of $[0.0, 1.2]$. The a priori profile for relative humidity
has rather arbitralily been chosen as

\begin{align}
\phi(t) = \begin{cases}
 0.7 &, 270\unit{K} < t \\
 0.7 - 0.01 \cdot (t - 270) & ,220 < t \leq  270\unit{K} \\
 0.2 \cdot (t - 270) & ,t < 220 \\
 \end{cases}.
\end{align}

The retrieval of liquid cloud mass density is restricted to the region
between the $0$ and $230\unit{^\circ C}$-isotherms. In contrast to the
frozen and liquid hydrometeors, cloud water is modeled in the retrieval
forward model to be purely absorbing using the absorption model by
\cite{liebe92} for suspended liquid cloud droplets. Liquid cloud mass density
is retrieved in $\text{log}_{10}$-space and the a priori profile is set to a fixed
value of $10^{-5}\ \unit{kg / m^3}$ in the permitted region of the atmosphere.

The a priori distributions of the 6 retrieval quantities ($N_0^*$ and $D_m$
for frozen and liquid hydrometeors, relative humidity $\phi$, cloud water)
are assumed to be independent so that the overall a priori covariance matrix
$\mathbf{S}_a$ has block-diagonal structure. Within each block, spatial
correlations between grid points are assumed to be exponentially decaying.
Hence, the correlation between grid points $i$ and $j$ of the values of the
retrieval quantity $q$ are computed as

\begin{align}
\left ( \mathbf{S}_{a,q} \right )_{i, j} &= \sigma_{q,i} \sigma_{q,j}
 \cdot \exp ( -\frac{d(i, j)}{l_q}),
\end{align}
where $\sigma_{q, i}$ is the a priori uncertainty assumed for retrieval
quantity $q$ at grid point $i$, $d(i, j)$ the distance between the grid
points and $l_q$ the quantity-specific correlation length. The assumed
a priori uncertainties and correlation lengths for the retrieval quantities
are summarized in Table \ref{tab:a_priori}.


\begin{table}[h!]
\caption{A priori uncertainties and correlation
 lengths used in the retrieval.}
 \centering
\label{tab:a_priori}
    \begin{tabular}{c|r|r}
     Quantity $q$ & $\sigma_q$ & $l_q$ [km]\\
    \hline
    $\log_{10}(N_{0, \text{frozen}}^*)$    & $2$                       & $2$ \\
    $D_{m, \text{ice}}$               & $300\unit{\mu m}$ & $2$ \\
    $\log_{10}(N_{0, \text{liquid}}^*)$    & $2                      $ & $2$ \\
    $D_{m, \text{liquid}}$            & $500\unit{\mu m}$ & $2$ \\
    $\text{arctanh}(\frac{2 \cdot \phi}{1.1} - 1.0)$ & $2$                       & $2$ \\
    $\log_{10}(m_\text{liquid cloud}) $ & $1$                       & $2$ \\
    \end{tabular}
\end{table}

As baselines for the assessment of the combined retrieval, also a radar-only and
a passive only-retrieval are performed. The radar-only retrieval uses the same
implementation as the combined retrieval, but only retrieves frozen and liquid
hydrometeors. For the radar-only retrieval, perfect knowledge of the atmospheric
humidity profile is assumed but liquid cloud is ignored in the retrieval forward
model.

The setup and retrieval quantities of the passive-only retrieval are similar to
the combined retrieval, with the only difference being that frozen and liquid
hydrometeors are retrieved at reduced resolution. For ice, $N_0^*$ is retrieved
at two equally spaced grid points between $0\unit{^\circ C}$-isotherm and
troposphere, while $D_m$ is retrieved at six. For liquid hydormeteors, the
retrieval grids for $N_0^*$ and $D_m$ are reduced to 2 equally spaced points
between surface and $0\unit{^\circ C}$-isotherm. Relative humidity is retrieved
at eleven grid points between surface and tropopause.

\section{Results}
\label{sec:results}

In the first part of this section, results of a numerical experiment are
presented that investigate the complementary information content of the active
and passive observations. The remainder of this section presents  the retrieval results
of the combined and the baseline retrievals applied to the reference cloud scenes.

\subsection{Complementary information content}
\label{sec:orgdc1b0b8}

A fundamental question regarding the benefit of combining two remote sensing
observations in a retrieval is to what extent the observations contain
non-redundant information. The degree of non-redundancy in the combined
observations is what is referred to here as complementary information content.
In order to explore this complementary information content in the radar and
radiometer observations, an idealized, homogeneous cloud layer of thickness of
$4\unit{km}$ located at an altitude of $10\unit{km}$ in a tropical atmosphere is
considered. The cloud is assumed to consist of a single frozen hydrometeor
species using the same PSD parametrization as the frozen hydrometeors used in
the cloud retrieval.

The question that we address here is whether and to what extent the active and
passive observations can constrain the two scaling factors of the PSD of frozen
hydrometeors in the cloud. To investigate this, the $N_0^*$ and $D_m$ parameters
of the PSD are varied and observations of the homogeneous cloud layer are
simulated. Figure \ref{fig:isolines} displays the simulated cloud signal, i.e.
the brightness temperature difference between clear sky and cloudy sky
simulation, as filled contours for a selection of channels of the MWI and ICI
sensors. For given values of $N_0^*$ and $D_m$, the ice mass density is given by
the relation
\begin{align}
m = \frac{\pi \rho}{4 ^ 4}N_0^* D_m^4.
\end{align}
Displaying the cloud signal in $D_m$-mass density space thus shows how the
measured passive cloud signal varies with the parameters of the PSD of the
frozen hydrometeors. Overlaid onto the contours of the passive cloud signal are
the isolines of the simulated radar reflectivity at the center of the cloud.

\begin{figure}
\centering
\includegraphics[width = 0.8\textwidth]{../plots/contours}
\caption{Simulated observations of a homogeneous cloud layer with
varying mass density $m$ and mass-weighted mean diameter $D_m$. The panels
display the radar reflectivity in dBZ at the cloud center overlaid onto the
cloud signal measured by selected radiometer channels of the MWI radiometer
(first row) and the ICI radiometer (second row).}
\label{fig:isolines}
\end{figure}

The contours of the measured active and passive cloud signals show the ambiguity
of each signal with respect to the two parameters of the PSD: Along these
contours the signal does not change, while the cloud composition does. A
necessary condition for the passive observations to provide complementary
information to the radar observations that can help to resolve the $N_0^*-D_m$
ambiguity is that the contours of the passive signal and the radar signal cross
each other. The panels in Fig.~\ref{fig:isolines} thus give an indication to
what extent the information in the radar measurement and the corresponding
passive radiometer channel provide complementary information on the two moments
of the PSD. Considering the panels corresponding to the MWI channels, the
results show that the observations contain complementary information only for
dense clouds consisting of large particle. Compared to that, the ICI
observations display a lower degree of parallelism already at lower densities
and particle sizes, indicating higher complementary information content already
for thinner clouds.

\subsection{Retrieval results}

To assess the performance of the combined cloud retrieval, the developed
algorithm has been applied to the two designated cloud scenes. The same
retrievals have been performed with a radar-only and a passive-only version of
the algorithm to serve as baselines for the combined retrieval. Each retrieval
was performed multiple times using different ice particle models. The tested
particle shapes are listed together with the corresponding mass size relations
and ARTS scattering database identifier in Table~\ref{tab:particles_retrieval}.
Since the results for both test scenes are qualitatively similar, not all
analyses are shown for both scenes. Instead, these are provided as a digital
supplement to this article.

\begin{table}
  \centering
  \caption{Particle model name, ARTS scattering database ID and parameters
    $\alpha, \beta$ of the mass-size relationships of the particle habits used
    in the retrieval.}
  \begin{tabular}{l|c|c|c}
    Name & ID & $\alpha$ & $\beta$ \\
    \hline
    GemCloudIce         & 31  & 440      & 3 \\
    GemSnow             & 32  & 24.0072  & 2.8571 \\
    GemGraupel          & 33  & 172.7527 & 2.9646 \\
    8-ColumnAggregate   &  8  & 65.4480  & 3      \\
    PlateType1          &  9  & 2.4770   & 0.7570 \\
    LargePlateAggregate &  20 & 2.2571   & 0.2085 \\
  \end{tabular}
  \label{tab:particles_retrieval}
\end{table}

The simulated observations for the first test scene are shown in
Figure~\ref{fig:observations_a}. To account for sensor noise, independent
Gaussian noise with standard deviations according to sensor specifications has
been added to the simulated observations. To model the forward model error that
would affect a real retrieval, the forward simulations have been performed using
the full microphysics from the GEM model. This means that the synthetic
observations were computed using six hydrometeor classes, while the retrieval
forward model uses only two classes of hydrometeors. Errors due mismatch of
sensor footprints or inhomogeneity of the atmosphere across the fields of view
of the sensors are not represented in the simulated observations.


\begin{figure}
\centering
\includegraphics[width = 0.8\textwidth]{../plots/observations_a}
\caption{Total hydrometeor content and simulated observations for the first
  test scene. Panel (a) displays the total hydrometeor content in the scene,
  i.e. the sum of the mass densities of all hydrometeor species of the GEM
  model. Panel (b) shows the simulated radar reflectivities for the scene.
  Panel (c) displays the simulated brightness temperatures for a selection
  of the channels of the MWI and ICI radiometers that are used in the retrieval.}
\label{fig:observations_a}
\end{figure}

\subsubsection{Mass concentrations}

To provide an overview over how well the different retrievals are able to
reproduce the cloud structures in the test scene, the retrieved ice mass
concentration fields for the first test scene are shown in Figure
\ref{fig:results_a}. The results shown in the figure were obtained using the
LargePlateAggregate as particle model. Panel (a) of the figure displays the
final value of the OEM cost function (normalized by the dimension of the
measurement space) achieved during the minimization process. Already here, a
first difference between the three retrievals becomes apparent: While the
radar-only retrieval achieves a low final cost over almost the whole scene, the
passive-only and the combined retrieval do not achieve a low cost in regions
where the cloud is thick and complex in structure.

Panel (c) shows the results of the passive-only retrieval. Although there are
similarities to the reference mass concentration, the results do not reproduce
the vertical structure of the cloud very well. It should be noted, however, that
the displayed mass-density range extends below the expected sensitivity limit of
the passive-only observations around $10\ \unit{mg/m^3}$, which explains the
smeared-out appearance of the results to some extent.

In contrast to this, the radar-only results, shown in panel (d), reproduce the
vertical structure of the cloud well. Nonetheless, when compared to the
reference mass concentration field, certain discrepancies are visible: The
radar-only retrieval tends to overestimate the mass density at the bottom
of the cloud and underestimate the mass concentrations at the top of the cloud.

The results of the combined retrieval are displayed in panel (e). Although some
artefacts are clearly visible in the retrieved mass density field, the retrieval
still reproduces the vertical structure well. In particular, the combined
retrieval succeeds to correct some of the systematic deviations of the
radar-only retrieval: The mass at cloud base is reduced and increased at cloud
top.


\begin{figure}
\centering
\includegraphics[width = 0.8\textwidth]{../plots/results_a_LargePlateAggregate}
\caption{Reference and retrieved mass concentrations of frozen hydrometeors for
 the first test scene. Panel (a) displays the final OEM costs achieved during
 the minimization normalized by the number of dimensions of the measurement
 space. Panel (b) displays the reference mass concentrations from the model
 scene. Panel (c), (d) and (e) display the retrieval results for the
 passive-only, radar-only and combined retrieval.}
\label{fig:results_a}
\end{figure}

To make the assessment of the retrieval performance more quantitiave, the
reference mass concentrations are plotted against the retrieved values in
Fig.~\ref{fig:results_scatter_a_1} and \ref{fig:results_scatter_a_2}. The plots
show the results for all different retrieval configurations and and tested
particle models. Markers in the plots are color-coded according to the
prevailing hydrometeor type (by mass density) in the reference scene, in order
to allow assessment of the retrieval performance for the different hydrometeor
types of the GEM model.

Not surprisingly, the results from the passive-only retrieval exhibit the
strongest deviations from the diagonal. Since the passive channels alone contain
only limited information on the vertical distribution of ice in the atmosphere,
no very good performance can be expected from the results at the vertical
resolution considered here. Although rather weak, a certain effect of the 
ice particle model on the retrieval results can be observed. In particular, the
GemCloudIce seems to lead to a systematic underestimation of ice mass densities,
that are less pronounced for the other particle models.

\begin{figure}
\centering
\includegraphics[width = 0.8\textwidth]{../plots/results_scatter_a_1}
\caption{Reference mass concentrations plotted against retrieved mass
  concentrations for the tested retrieval configurations. Each row shows the
  retrieval results for the particle shape shown in the first panel of each row.
  The following panels show the retrieval results for the passive only (second
  column), the radar only (third column) and the combined retrieval (fourth
  column). Markers are colored according to the prevailing hydrometeor type at
  the corresponding grid point in the reference scene. Due to their sparsity,
  marker corresponding to hail are drawn at twice the size of the other
  markers.}
\label{fig:results_scatter_a_1}
\end{figure}

The results from the radar-only retrieval are more accurate than the passive-only
retrieval, with almost all retrieval results located fairly close to the diagonal.
The most distinct feature of the radar-only results, however, is the emergence
of two clusters extending along the diagonal displaced above and below it. The
color coding of the markers reveals that these clusters correspond to grid
points dominated by ice for the clusters below the diagonal and snow for the
cluster located above the diagonal. This indicates that the radar-only retrieval
systematically underestimates mass densities for cloud ice but overestimates the
mass density of snow. The effect is observed for all tested particle shapes
indicating that it is probably related to the microphysical assumptions
regarding the PSD rather than particle shape. In general, the radar-only results
exhibit only very weak dependency on the assumed particle shape making the
radar-only results for different assumed particle shapes virtually
indistinguishable.

Another feature that stands out in the radar-only results, is that the retrieval
does not work for graupel. This, however, can be understood by comparing the
radar reflectivities shown in Fig.~\ref{fig:observations_a} with the cloud
structure displayed in Fig.~\ref{fig:overview}. It becomes apparent that graupel
in this scene is located where the radar signal is fully attenuated. Since there
is no signal to retrieve the mass density from, this explains the bad
performance of the radar-only retrieval for these grid points.

Similar as the radar-only retrieval, the results of the combined retrieval are
located close to the diagonal. However, the clusters observed in the radar-only
result are to large extent merged in the combined results. Moreover, except for
the results obtained with the GemCloudIce particle shape, the two clusters move
in closer towards the diagonal. This improved retrieval performance of the
combined retrieval, clearly indicates that the complementary information content in
the passive observations helps to better constrain the cloud microphysics in the
different regions of the cloud.

Nonetheless, the results for the GemCloudIce particle stand out in the results.
Even though the systematic deviations observed in the radar-only retrieval are
reduced for most particle shapes, for this specific shape they are instead
increased. This is caused mostly by the effect the particle shape has on the
retrieval of snow, which is strongly underestimated for reference mass
concentrations around $10^{-4} \ \unit{kg/m^3}.

\begin{figure}
\centering
\includegraphics[width = 0.8\textwidth]{../plots/results_scatter_a_2}
\caption{Same as Fig.~\ref{fig:results_scatter_a_1} but for the remaining particle
  shapes.}
\label{fig:results_scatter_a_2}
\end{figure}


The results for the second test scene obtained using the LargePlateAggregate
particle model are shown in Fig.~\ref{fig:results_b}. As mentioned above, the
results are qualitatively very similar to those of the first scene. Also here,
the final OEM cost, shown in Panel (a), displays a region of increased cost for
the passive-only and combined retrievals. This is again a region of very dense
cloud that consists of graupel and snow. Also similar to the first scene, the
passive only retrieval does not reproduce the structure of the cloud well.
Although the cloud top is placed at the right position neither the vertical
structure of the cloud nor its base are resolved. The radar-only retrieval
resolves the vertical structure of the cloud well, but overestimates the ice
mass density in the scene. The combined retrieval also resolves the vertical
structure of the cloud well and corrects the overestimation observed in the
radar-only results to some extent.

\begin{figure}
\centering
\includegraphics[width = 0.8\textwidth]{../plots/results_b_LargePlateAggregate}
\caption{Reference and retrieved mass concentrations of frozen hydrometeors for
 the second test scene. Panel (a) displays the final OEM costs achieved during
 the minimization normalized by dimensions of the measurement space. Panel (b)
 displays the reference mass concentrations from the model scene. Panel (c),
 (d) and (e) display the retrieval results for the passive-only, radar-only
 and combined retrieval.}
\label{fig:results_b}
\end{figure}

Scatter plots for the retrieval results from the second scene are shown in
Fig.~\ref{fig:results_scatter_b_1}. Except for the lack of cloud ice in the
scene, the results are the same as what has been observed in the first scene:
The radar-only retrieval overestimates the mass density of snow in the scene.
This effect is corrected by the combined retrieval for most of the tested
particle shapes. The exception is the GemCloudIce particle for which the
retrieval of snow particle deteriorates quite drastically.

\begin{figure}
\centering
\includegraphics[width = 0.8\textwidth]{../plots/results_scatter_b_1}
\caption{Scatter plots of the reference and retrieved ice mass densities for
  the second test scene. The rows show the retrieval results for a given
  assumed ice particle model. The first column of each row displays a rendering
  of the particle model. The following rows display the results for the
  passive-only, the radar-only and the combined retrieval.}
\label{fig:results_scatter_b_1}
\end{figure}


\subsubsection{Particle number densities}

A further interesting question is whether and to what extent the joint
observations can improve the retrieval of particle number concentrations.
To address it, particle number densities have been computed from the
retrieved parameters by computing the zeroth moment of the PSD. The resulting
particle number density fields are displayed together with the reference
field in Fig.~\ref{fig:results_nd_a}.

Comparing the passive-only and the radar-only retrieval to the reference
field, the results show that both methods lack skill in predicting number
density concentrations. The strong gradient between the very high concentrations
at the top of the cloud and the low concentrations at the bottom is not resolved
by either of the methods. The combined retrieval, however, manages to reproduce
this gradient in most parts of the scene. Although its exact structure
is not fully reproduced, this clearly shows sensitivity of the retrievals to
the number density concentrations.

Of interest is also the part of the scene where the combined retrieval shows
the strongest deviations from the reference field, that is between $2$ and
$3\unit{^\circ}$ latitude. Here, the results strongly underestimate the true
number concentrations. Comparison with the cloud composition displayed Panel (a)
of Fig.~\ref{fig:overview} shows that this region matches the regions of the
cloud, that contain large amounts of both cloud ice and snow. Since the
retrieval uses only a single hydrometeor species to represent ice in the
atmosphere, it is not able to represent such heterogeneous conditions. Since
snow will have the stronger impact on the observations, the retrieval in these
regions tends to predict snow rather than ice, which leads to the low retrieved
number densities.

\begin{figure}
\centering
\includegraphics[width = 0.7\textwidth]{../plots/results_nd_a_LargePlateAggregate}
\caption{Reference and retrieved particle number concentrations of frozen
  hydrometeors for the first test scene obtained with the LargePlateAggregate
  particle model. Panel (a) displays the reference mass concentrations from the
  model scene. Panel (b), (c) and (d) display the retrieval results for the
  passive-only, radar-only and combined retrieval.}
\label{fig:results_nd_a}
\end{figure}

Figure~\ref{fig:results_nd_scatter_a} displays scatter plots of the reference
and retrieved number density concentrations for all three methods and two
particle models from the first test scene. Markers in the plot are color coded
according to a purity index, which has been defined as the ratio of the maximum
mass density of any of the frozen hydrometeor species in the model and the total
mass density. The color coding reveals the emergence of two clusters in the
results: One located at low reference particle number concentrations
corresponding to snow and one at high number concentrations corresponding to
regions where the cloud consists of cloud ice or graupel.

The results confirm that the radar-only retrieval does not exhibit any skill in
retrieving particle number densities. The passive-only results, however, seem to
indicate at least some sensitivity since the cluster corresponding to snow is
placed correctly on the diagonal. In contrast to that, the combined retrieval
moves both clusters towards the diagonal, indicating that it is capable of
distinguishing the microphysical differences of cloud ice and snow. Yet still,
the accuracy of the combined retrieval remains very low.

The effect of particle shape on the retrieval results is similar to what has
been observed for the mass density retrievals. The GemCloudIce model yields the
worst retrieval results, leading to a general underestimation of the true
particle number density. Results for the other particles are mostly similar,
which is why only the result for the LargePlateAggregate is shown here.

\begin{figure}
\centering
\includegraphics[width = 0.7\textwidth]{../plots/results_nd_scatter_a}
\caption{Scatter plots of the retrieved particle number densities at grid points with
  reference mass density larger than $10^{-5}\ \unit{kg/m^3}$. Rows show the results
  for the different particle models used in the retrieval while column display the
  results for the different retrieval methods. Marker color encodes a purity index 
  which grid point which is computed as the ratio of the maximum mass density of
  any of the frozen hydrometeor species and the total mass density.}
\label{fig:results_nd_scatter_a}
\end{figure}

\subsubsection{Impact of assumed ice particle shape}

To better understand the effect that the assumed ice particle shape has on the
retrieval results, the mass density relations for the tested particle models are
displayed in panel (a) of Fig.~\ref{fig:costs}. As can be seen from this plot,
the mass size relationship of the GemCloudIce particle clearly stands out as
being much denser than the others. On the contrary, the other particles have
quite similar mass size relationships. This seems to explain why the GemCloudIce
particle model does not work well for snow in the combined retrievals,
explaining the observed deviations in the retrieved mass densities.

Also displayed in Fig~\ref{fig:costs} (panel (b) and (c)) are the final OEM
costs of the combined retrieval that were obtained with the different particle
models for both test scenes. Since the particle shape affects the sub-millimeter
observations, one could hope that the retrieval results can be used to infer the
prevailing ice particle shape within the cloud based on the quality of the fit
of the forward model achieved during the retrieval.

Unfortunately, such clear conclusions cannot be drawn from the results. In the
first test scene, the best fit is obtained by the GemSnow, GemCloudIce and the
LargePlateAggregate particle models, although the GemCloudIce model clearly
resulted in the worst retrieval performance. For the second scene, on the other
hand, there appears to be better agreement between the lowest retrieval cost and
retrieval performance. Here, the GemSnow particle consistently gives the lowest
OEM cost which is also the dominant frozen hydrometeor species in the scene.

\begin{figure}
\centering
\includegraphics[width = 0.4\textwidth]{../plots/costs}
\caption{Mass-size relations (Panel (a)) and final OEM cost for the two test scenes
  (Panel (b) and (c)). The final cost curves where smoothed using a running average filter
of a width of 20 profiles.}
\label{fig:costs}
\end{figure}


\subsubsection{Humidity and cloud water}

The developed passive and combined retrieval algorithms also retrieve profiles
of humidity and liquid cloud mass density. For relative humidity, both
retrievals showed a certain sensitivity but no improvement could be observed in
the results of the combined retrieval compared to the passive-only retrieval.
Moreover, no suitable retrieval setup was found that would achieve satisfactory
performance within the scope of this study. Considering that the sensor setup
with only channels with frequencies at $89\ \unit{GHz}$ and higher is not ideal for
humidity retrievals, we do not consider our results representative of what could
be achieved with the observational approach and thus do not include them here.

Nonetheless, we provide results from the liquid cloud retrieval in
Fig.~\ref{fig:results_cw_b} since these show another interesting synergy of the
radar and passive microwave observations. The results in
Fig.~\ref{fig:results_cw_b} display contours of the retrieved mass density field
of liquid cloud drawn on top of the retrieved frozen hydrometeor mass density
field for the passive-only and the combined retrieval. These results show
clearly that the combined retrieval is able to detect and retrieve liquid clouds
even when they overlap with ice clouds. Although some sensitivity of the
retrieval to liquid clouds can be observed also in the passive retrieval, the
cloud is put at the wrong altitude and its mass density is greatly
underestimated. This indicates that the radar reflectivity profile contains
useful information for the retrieval to better locate cloud water in the
atmospheric column.

\begin{figure}
\centering
\includegraphics[width = \textwidth]{../plots/results_cw_b_LargePlateAggregate}
\caption{Contours of the cloud water mass density field drawn on top of the retrieved
  frozen hydrometeor mass density field. The results shown were obtained using the
  LargePlateAggregate particle model.}
\label{fig:results_cw_b}
\end{figure}


\section{Discussion}
\label{sec:discussion}

The principal aim of this study was to investigate the synergies between radar
and passive sub-millimeter observations. To this end, a simplified numerical
experiment has been presented, that clearly demonstrates the complementarity of
the information contained in the radar and the sub-millimeter observations.
Furthermore, a combined retrieval algorithm has been developed to demonstrate
the feasibility of the synergistic retrievals and further explore their
potential and limitations.

\subsection{Fundamental synergies}

The experiment presented in the first part of this study aimed to establish the
fundamental synergies of the active and passive microwave observations. It
compared the cloud signals observed by a radar, a millimeter-wave radiometer and
a sub-millimeter-wave radiometer. The results show that the combined
observations can constrain the horizontal as well as the vertical scaling of the
particle size distribution. Furthermore, the complementary information content
between the active and passive observations was shown to depend on both the
properties of the observed cloud and the frequency of the observations. For the
lower frequencies considered in this study, i.e. the highest channels of the MWI
radiometer, the regions where both observations provide complementary
information on the particle size distribution of the cloud are limited to high
mass densities and particle sizes. As the passive observing frequency increases,
the regions of complementary information content extend down to smaller particle
sizes and cloud mass density. Especially the highest-frequency channels of the
ICI radiometer can therefore be expected to provide additional information on
the particle size distribution of ice clouds.

These results indicate that the combination of active and passive microwave
observations provides different synergies depending on the frequency of the
passive observations: While the MWI channels will help to better constrain the
distributions of precipitating hydrometeors, improvements of the retrieval of
cloud ice can only be expected from the sub-millimeter observations.

\subsection{Combined cloud retrieval}

In the second part of the study, we have presented results from a combined,
variational cloud retrieval applied to synthetic observations from two test
scenes from a high-resolution atmosphere model. The results of the combined
retrieval were compared to that of a passive- and a radar-only version of the
retrieval algorithm. The simulated observations neglected potential errors
caused by different or non-overlapping antenna beams as well as inhomgeneity of
the atmosphere across the beams. A source of forward model error was included in
the simulated observations by performing them applying the full microphysics
scheme from the GEM model. This allowed to assess of the effect of cloud
composition as well the simplified microphysics used in the retrieval forward
model on the results.

Of the three considered retrieval implementations, the passive-only retrieval
clearly yields the worst performance for the retrieval of mass-density profiles.
It should be noted, however, that the passive only retrieval presented here has
not been optimized to its full potential and should therefore not be taken as
representative for the performance of the MWI and ICI radiometers. To ensure a
fair comparison, the passive only retrieval uses very similar a priori
assumptions as the other two retrievals, which in the presented case provide
only very limited information on the vertical structure of the cloud. As other
studies have shown, the passive observations do provide information on the
vertical distribution of ice in the atmospheric column \citep{wang16,
  grutzun18}, but the information content is limited to a few degrees of
freedom. It is therefore unlikely, the vertical resolution of the passive-only
can be improved drastically without further constraining it a priori, for
example through an observation database in an empirical retrieval.

As expected, the radar-only retrieval provides much better profile retrievals
than the passive-only version. However, the results exhibit systematic
deviations from the reference mass densities. Closer inspection revealed that
these are caused by an underestimation of the mass density of cloud ice and a
simultaneous overestimation of the mass density of snow. This indicates that the
a priori assumptions used in the retrieval do not fit the way ice and snow are
represented in the model, which causes the systematic deviations in the
retrieved mass densities. This observation is confirmed by the number density
fields retrieved from the radar-only retrieval. These are too low at regions
consisting mostly of cloud ice and too high in regions consisting of mostly
snow. Since the radar-only number density retrieval does not show signs of
sensitivity to the particle number density in the cloud, the reason for this
retrieval error are likely the a priori-assumed microphysical properties of the
hydrometeors.

The a priori assumptions that were used in this study were similar but not
identical to what is used in the DARDAR retrievals. Also here it should be
noted, that the presented results should not be taken to be representative for
the DARDAR product. Rather than this, the DARDAR a priori settings were chosen
since they represent well established and validated assumptions for ice cloud
retrievals and therefore should provide a reasonable starting point for the
development of a combined cloud retrieval. The fact that the a priori
assumptions used in the DARDAR retrieval do not agree with the microphysical
properties of ice and snow in the model, does not say much about the general
validity of these assumptions.

Despite the fact that the combined retrieval produced certain artefacts in the
results, our analysis showed that it was, at least for most of the tested
particle models, able to correct the systematic deviations in the retrieved mass
density field, that were observed in the radar-only retrieval. The benefit of
the combined observations was even more pronounced in the retrieved number
density fields. Here, neither the passive- or radar-only retrieval showed any
skill in retrieving the particle number concentrations. The combined retrieval,
however, was able to reproduce the general structure of the number concentration
fields in regions where the cloud composition was homogeneous. Despite the
numerical accuracy of the number density is still low, this clearly shows that
the joint observations constrain the cloud microphysics better than the passive
or radar observations alone.

Although the presented results show that the combined observations can improve
the ice cloud retrieval, the overall performance of the combined retrieval can
even degrade if performed with an unrealistic particle habit. Moreover, the
combined retrieval displayed a stronger dependence of the retrieval error on the
assumed particle shape than the radar-only retrieval did. The higher sensitivity
of the sub-millimeter channels thus increases the importance of a realistic
particle model. Unfortunately, the final OEM cost function value did not display
any clear relation to most the suitable particle model used in the retrieval.
While this is a negative results for this study, it certainly does not mean that
the observations do not provide any information on the particle shape. A good
fit of a particle model in certain parts of the atmosphere can easily be swamped
by other forward model errors or misfits in other parts of the atmospheric
column, etc.

As an outlook, we have also included results of a liquid cloud retrieval, that
clearly shows the capability of the combined retrieval to distinguish liquid
clouds from ice clouds and retrieve profiles of liquid cloud mass densities.
Since the passive-retrieval alone was not able to retrieve the liquid cloud,
this indicates an additional synergy between the active and the passive
observations. Although no satisfactory results were obtained from the water
vapor retrieval, the retrieval results (not shown) still indicate sensitivity of
this setup for retrieving atmospheric humidity. The full exploration of the
potential of the combined observations for liquid cloud and water vapor is out
of the scope of this study and is left to future investigation.

The combined retrieval implementation showed robust performance on fairly
distinct and complex cloud scenes. Despite this, both scenes that were
considered here contained parts where the OEM minimization did find a state that
results in a good fit to the observations. In contrast to that, the radar-only
retrieval did achieve a was low in most regions where the final cost of the
combined retrieval was high. This indicates that there is actually even more
information contained in the combined observations, that the retrieval is not
able to disentangle. This raises the question of the suitability of the OEM
method applied here. The combined retrieval violates the two fundamental
assumptions of the OEM method: The forward model is strongly non-linear and the
assumed Gaussian a priori assumptions are likely not a very good description of
reality. In addition to that, the current implementation of the retrieval is
computationally too expensive to be applied operationally. For further
development of the combined retrieval concept it may therefore be advisable to
revisit the applied retrieval method in search for a potentially more suitable
alternative.

Finally, it is important to consider the limitations of this study. Its primary
limitation is certainly the restriction to the model scenes for the testing of
the retrieval. In addition to that, the forward simulations used to generate the
synthetic observations do not consider beam filling issues, assume a slightly
unrealistic viewing geometry and neglect multiple scattering in the radar
simulations. It is therefore important to understand the results presented here
in their relative sense rather than providing realistic performance
characteristics of the considered retrieval methods. What the presented results
show is that, due to the synergies of the active and passive observations, the
combined retrieval is able to better predict the microphysical properties of the
observed clouds than any of the single-instrument retrievals alone.

\conclusions  %% \conclusions[modified heading if necessary]
\label{sec:conclusions}

\begin{itemize}
  - Demonstrated existence of synergies between the active and passive
  microwave observations
  - Retrieval of liquid clouds as an additional synergy
  - Mass density as well as number density better constrained in combined
  retrieval
  - Particle shape affects retrieval, but no clear best fit was found
\end{itemize}

The results of this study clearly prove the existence of synergies between radar
and sub-millimeter observations. For the retrieval of frozen hydrometeors,
combining active with the passive radiometer observations helps to improve both
the retrieval of mass densities as well as particle number concentrations.
From these results together with the conceptual considerations based on the
analysis of the cloud signal in the active and passive observations it can be
concluded that the combined observations help to better constrain the
microphysical properties of the observed ice particles.

The retrieval results also point towards an additional synergy between the
acitve and passive observations: The detection and retrieval of liquid clouds.
Neither of the active-only or passive-only retrievals,  was capable
In addition to this, the presented retrieval results also show potential for retrieval of liquid clouds. SInc e


  
The results presented in this study clearly show the existence of synergies
between sub-millimeter radiometer and radar observations for the retrieval of
frozen hydrometeors. The combined observations allow the retrieval of two
moments of the particle size distribution and thus help to better constrain
both the mass density as well as the particle number density in the cloud.
In addition to that, our results indicate additional synergies for the retrieval
of liquid clouds and water vapor.

In this study, we have investigated fundamental synergies and demonstrated the
feasibility of the combined retrievals. This work hence opens up a number of
opportunities for further research: In addition to the rather obvious
application of the retrieval to real observations, there is the exploration of
the potential of the combined observations for the retrieval of humidity and
liquid cloud. Fruthermore, our results have shown that a realistic ice particle
model is required for the retrieval to work well. This highlights the need for a
realistic particle model that is suitable for wide range of cloud configurations
encountered in the real atmosphere. Finally, another question to address is
whether there are suitable alternatives to the OEM method to use as a retrieval
method for the combined retrieval.

Since only microwave observations can provide sensitivity to frozen hydrometeors
throughout the full atmospheric column in the presence of thick clouds, the
concept of combining active and passive microwave observations has great
potential to reduce uncertainties on global observations of ice in the
atmosphere. Although this work constitutes only a small step towards
fully-functional retrievals from combined radar and sub-millimeter observations,
it clearly shows the potential of the approach. The ability of the joint
observations to better constrain the microphysics of clouds should make them an
interesting concept to consider for future field campaigns or space missions.


%% The following commands are for the statements about the availability of data sets and/or software code corresponding to the manuscript.
%% It is strongly recommended to make use of these sections in case data sets and/or software code have been part of your research the article is based on.

\codeavailability{All code used to produce the results in this study is available through a public repository \citep{ccrf}.} %% use this section when having only software code available


\dataavailability{} %% use this section when having only data sets available



\appendix
\section{}    %% Appendix A

\subsection{}     %% Appendix A1, A2, etc.


\noappendix       %% use this to mark the end of the appendix section

%% Regarding figures and tables in appendices, the following two options are possible depending on your general handling of figures and tables in the manuscript environment:

%% Option 1: If you sorted all figures and tables into the sections of the text, please also sort the appendix figures and appendix tables into the respective appendix sections.
%% They will be correctly named automatically.

%% Option 2: If you put all figures after the reference list, please insert appendix tables and figures after the normal tables and figures.
%% To rename them correctly to A1, A2, etc., please add the following commands in front of them:

\appendixfigures  %% needs to be added in front of appendix figures

\appendixtables   %% needs to be added in front of appendix tables

%% Please add \clearpage between each table and/or figure. Further guidelines on figures and tables can be found below.



\authorcontribution{Simon Pfreundschuh has implemented the retrieval, performed the data analysis and written
  the manuscript. Patrick Eriksson and Richard Larsson have contributed code to the ARTS radiative transfer
  model that was required to perform the presented calculations. Stefan A. Buehler, Patrick Eriksson, Manfred
  Brath and Simon Pfreundschuh have collaborated on the study that lead to the results presented here. David Duncan
and Robin Ekelund have contributed to the conceptualization of the study through comments and advice.} %% this section is mandatory for the journals ACP and GMD. For all other journals it is strongly recommended to make use of this section

\competinginterests{No competing interests are present} %% this section is mandatory even if you declare that no competing interests are present

\begin{acknowledgements}
TEXT
\end{acknowledgements}




%% REFERENCES

%% The reference list is compiled as follows:


%% Since the Copernicus LaTeX package includes the BibTeX style file copernicus.bst,
%% authors experienced with BibTeX only have to include the following two lines:

\bibliographystyle{copernicus}
\bibliography{references.bib}
%%
%% URLs and DOIs can be entered in your BibTeX file as:
%%
%% URL = {http://www.xyz.org/~jones/idx_g.htm}
%% DOI = {10.5194/xyz}


%% LITERATURE CITATIONS
%%
%% command                        & example result
%% \citet{jones90}|               & Jones et al. (1990)
%% \citep{jones90}|               & (Jones et al., 1990)
%% \citep{jones90,jones93}|       & (Jones et al., 1990, 1993)
%% \citep[p.~32]{jones90}|        & (Jones et al., 1990, p.~32)
%% \citep[e.g.,][]{jones90}|      & (e.g., Jones et al., 1990)
%% \citep[e.g.,][p.~32]{jones90}| & (e.g., Jones et al., 1990, p.~32)
%% \citeauthor{jones90}|          & Jones et al.
%% \citeyear{jones90}|            & 1990



%% FIGURES

%% When figures and tables are placed at the end of the MS (article in one-column style), please add \clearpage
%% between bibliography and first table and/or figure as well as between each table and/or figure.


%% ONE-COLUMN FIGURES

%%f
%\begin{figure}[t]
%\includegraphics[width=8.3cm]{FILE NAME}
%\caption{TEXT}
%\end{figure}
%
%%% TWO-COLUMN FIGURES
%
%%f
%\begin{figure*}[t]
%\includegraphics[width=12cm]{FILE NAME}
%\caption{TEXT}
%\end{figure*}
%
%
%%% TABLES
%%%
%%% The different columns must be seperated with a & command and should
%%% end with \\ to identify the column brake.
%
%%% ONE-COLUMN TABLE
%
%%t
%\begin{table}[t]
%\caption{TEXT}
%\begin{tabular}{column = lcr}
%\tophline
%
%\middlehline
%
%\bottomhline
%\end{tabular}
%\belowtable{} % Table Footnotes
%\end{table}
%
%%% TWO-COLUMN TABLE
%
%%t
%\begin{table*}[t]
%\caption{TEXT}
%\begin{tabular}{column = lcr}
%\tophline
%
%\middlehline
%
%\bottomhline
%\end{tabular}
%\belowtable{} % Table Footnotes
%\end{table*}
%
%%% LANDSCAPE TABLE
%
%%t
%\begin{sidewaystable*}[t]
%\caption{TEXT}
%\begin{tabular}{column = lcr}
%\tophline
%
%\middlehline
%
%\bottomhline
%\end{tabular}
%\belowtable{} % Table Footnotes
%\end{sidewaystable*}
%
%
%%% MATHEMATICAL EXPRESSIONS
%
%%% All papers typeset by Copernicus Publications follow the math typesetting regulations
%%% given by the IUPAC Green Book (IUPAC: Quantities, Units and Symbols in Physical Chemistry,
%%% 2nd Edn., Blackwell Science, available at: http://old.iupac.org/publications/books/gbook/green_book_2ed.pdf, 1993).
%%%
%%% Physical quantities/variables are typeset in italic font (t for time, T for Temperature)
%%% Indices which are not defined are typeset in italic font (x, y, z, a, b, c)
%%% Items/objects which are defined are typeset in roman font (Car A, Car B)
%%% Descriptions/specifications which are defined by itself are typeset in roman font (abs, rel, ref, tot, net, ice)
%%% Abbreviations from 2 letters are typeset in roman font (RH, LAI)
%%% Vectors are identified in bold italic font using \vec{x}
%%% Matrices are identified in bold roman font
%%% Multiplication signs are typeset using the LaTeX commands \times (for vector products, grids, and exponential notations) or \cdot
%%% The character * should not be applied as mutliplication sign
%
%
%%% EQUATIONS
%
%%% Single-row equation
%
%\begin{equation}
%
%\end{equation}
%
%%% Multiline equation
%
%\begin{align}
%& 3 + 5 = 8\\
%& 3 + 5 = 8\\
%& 3 + 5 = 8
%\end{align}
%
%
%%% MATRICES
%
%\begin{matrix}
%x & y & z\\
%x & y & z\\
%x & y & z\\
%\end{matrix}
%
%
%%% ALGORITHM
%
%\begin{algorithm}
%\caption{...}
%\label{a1}
%\begin{algorithmic}
%...
%\end{algorithmic}
%\end{algorithm}
%
%
%%% CHEMICAL FORMULAS AND REACTIONS
%
%%% For formulas embedded in the text, please use \chem{}
%
%%% The reaction environment creates labels including the letter R, i.e. (R1), (R2), etc.
%
%\begin{reaction}
%%% \rightarrow should be used for normal (one-way) chemical reactions
%%% \rightleftharpoons should be used for equilibria
%%% \leftrightarrow should be used for resonance structures
%\end{reaction}
%
%
%%% PHYSICAL UNITS
%%%
%%% Please use \unit{} and apply the exponential notation


\end{document}
